% Research contribution for the parent manuscript; not a standalone document.
% Suggested citation key: Panzer2017 (arXiv:1512.04482v3).

\section{Proof of the five Gaussian weight-six identities}
\label{sec:gaussian-six-proof}

Put
\[
 D_{a,b}(z)=\operatorname{Li}_{a,b}(z,1)
   =\sum_{n>m\ge1}\frac{z^n}{n^a m^b},\qquad
 g_{a,b}=\Im D_{a,b}(i),\qquad w=a+b.
\]
We use the principal analytic continuation from the unit disk. In particular,
on the open upper unit semicircle, $D_{a,b}(1/z)=\overline{D_{a,b}(z)}$.
The case $a=1$ is interpreted by the convergent boundary series, or
equivalently its Abel limit. For $b>1$, this boundary convergence follows
by writing $H_{n-1}^{(b)}=\zeta(b)+O(n^{1-b})$; for $b=1$, Dirichlet's
test applies to the eventually decreasing sequence $H_{n-1}/n$.

The five identities labelled \texttt{gauss:eq:g51} through
\texttt{gauss:eq:g15} in the source manuscript are all exact. Their
coefficients require no correction. What changes is their status: the
finite formulas below give proofs and reproducible rational certificates.
The general parity principle and its depth-two formula are known; the
contribution here is their careful specialization to the manuscript's
convention, explicit verification of the five rows, and a complete extension
to weight eight. No claim of priority for the general parity theorem is made.

\subsection{A finite inversion formula with the boundary cancellation exposed}

For $\Im z>0$, write $L=\Log z$ and set
\[
 \mathcal B_n(z)=\frac{(2\pi i)^n}{n!}
 B_n\!\left(\frac{L}{2\pi i}\right),\qquad n\ge0,
\]
where $B_1(x)=x-\tfrac12$. Thus $\mathcal B_0=1$ and
$\mathcal B_n(i)=(2\pi i)^nB_n(1/4)/n!$.

\begin{theorem}[Explicit depth-two inversion at $(z,1)$]
\label{thm:gaussian-inversion}
For integers $a,b\ge1$ and $\Im z>0$,
\begin{align}
 D_{a,b}(z)-(-1)^wD_{a,b}(1/z)
 ={}&-\operatorname{Li}_w(z)\label{eq:gaussian-inversion}\\
 &+\sum_{k=1}^{a}(-1)^k
   \binom{b+k-1}{b-1}\zeta(b+k)\mathcal B_{a-k}(z)\notag\\
 &+(-1)^a\sum_{k=0}^{b}\binom{a+k-1}{a-1}
   \operatorname{Li}_{a+k}(1/z)\mathcal B_{b-k}(z).\notag
\end{align}
In particular, no divergent $\zeta(1)$ occurs, including when $a=1$ or
$b=1$.
\end{theorem}

\begin{proof}
Apply the depth-two functional equation of Panzer,
\cite[Eq.~(3.2)]{Panzer2017}, with his ascending-index pair $(n_1,n_2)=(b,a)$
and arguments $(z_1,z_2)=(u,z)$, taking $u=e^{i\varepsilon}$ with
$\varepsilon\downarrow0$ small enough that $uz$ remains in the upper
half-plane. His Bernoulli factor
is exactly $\mathcal B_n(z)$ on the indicated branch, since
$\Log(-z)=L-i\pi$ in the upper half-plane. The term with $\mu=b$ in
the first sum and the separate product term cancel: their difference is
\[
 \operatorname{Li}_b(u)
 \bigl(\mathcal B_a(uz)-\mathcal B_a(z)\bigr)\longrightarrow0.
\]
For $b\ge2$ this is continuity. For $b=1$, the first factor is
$O(|\log(1-u)|)$ and the second $O(|1-u|)$, so it still tends to zero.
All remaining zeta arguments are at least $b+1\ge2$.
Reindexing the surviving sums gives \eqref{eq:gaussian-inversion}.
The left side has the required limit because both relevant multiple
polylogarithms are analytic near $(1,z)$ when $z$ is nonreal. The
independent differential proof below also verifies the specialization
without an appeal to a formal divergent substitution.
\end{proof}

\subsection{A rational coefficient formula at every even weight}

Let
\[
 q_n=\frac{2^n}{n!}B_n(1/4),\qquad
 C_{\mu,r}=\begin{cases}\binom{\mu-1}{r-1},&\mu\ge r,\\0,&\mu<r,
 \end{cases}
 \qquad d_\mu=2^{-\mu}(1-2^{1-\mu})\quad(\mu\ge2).
\]
Every $q_n$ is rational. Recall $\beta(s)=\sum_{j\ge0}(-1)^j/(2j+1)^s$
and $G=\beta(2)$.

\begin{corollary}[Explicit Gaussian coefficient rule]
\label{cor:gaussian-coefficient-rule}
If $w=a+b$ is even, then
\begin{align}
 g_{a,b}={}&\frac12\sum_{\substack{2\le\mu\le w\\\mu\ \mathrm{even}}}
 \left[(-1)^{a+1+(w-\mu)/2}C_{\mu,a}q_{w-\mu}
       -\delta_{\mu,w}\right]\pi^{w-\mu}\beta(\mu)
 \label{eq:gaussian-real-coefficients}\\
 &+\frac12\sum_{\substack{3\le\mu<w\\\mu\ \mathrm{odd}}}
 (-1)^{a+1+(w-\mu-1)/2}q_{w-\mu}
 \left[C_{\mu,b}-\delta_{\mu,b}+d_\mu C_{\mu,a}\right]
 \pi^{w-\mu}\zeta(\mu)\notag\\
 &+\frac{\delta_{a,1}}4(-1)^{w/2-1}q_{w-1}\pi^{w-1}\log2.\notag
\end{align}
Consequently $g_{a,b}$ lies in the rational linear span of
$\pi^{w-2j}\beta(2j)$, $\pi^{w-2j-1}\zeta(2j+1)$, and, only if $a=1$,
$\pi^{w-1}\log2$. This is a linear reduction to single special values
times powers of $\pi$.
\end{corollary}

\begin{proof}
At $z=i$ the left side of \eqref{eq:gaussian-inversion} is $2ig_{a,b}$.
Insert
\[
 \operatorname{Li}_1(-i)=-\tfrac12\log2-\tfrac{i\pi}{4},\qquad
 \operatorname{Li}_\mu(-i)=-d_\mu\zeta(\mu)-i\beta(\mu)\quad(\mu\ge2),
 \qquad\mathcal B_n(i)=i^n\pi^nq_n.
\]
Even $\mu$ in the last sum supplies its beta term; odd $\mu$ supplies
its real single-polylogarithm part. The first sum contributes only at odd
$\mu=b+k$, with $\mu>b$. This last strict inequality is represented by
$C_{\mu,b}-\delta_{\mu,b}$. Taking imaginary parts and dividing by two
gives the displayed formula.
\end{proof}

Two useful checks follow immediately. The coefficient of $\beta(w)$ is
\[
 \frac{(-1)^{a+1}\binom{w-1}{a-1}-1}{2},
\]
so it vanishes in $g_{1,w-1}$. Also $\beta(2j)$ can occur only when
$2j\ge a$, apart from the already included collision term at $2j=w$.
These structural zeros explain several sparse rows in the tables; they
are not conclusions drawn from a numerical search.

\subsection{The five manuscript rows are theorems}

Applying \eqref{eq:gaussian-real-coefficients} at $w=6$ gives
\begin{align}
 2048g_{5,1}&=-64\pi^3\zeta(3)-527\pi\zeta(5)+4096\beta(6),\\
 1536g_{4,2}&=96\pi^3\zeta(3)-32\pi^2\beta(4)
              +1581\pi\zeta(5)-8448\beta(6),\\
 1024g_{3,3}&=-3\pi^3\zeta(3)+64\pi^2\beta(4)
              -1581\pi\zeta(5)+4608\beta(6),\\
 23040g_{2,4}&=-14\pi^4G+135\pi^3\zeta(3)-1440\pi^2\beta(4)
               +23715\pi\zeta(5)-69120\beta(6),\\
 92160g_{1,5}&=-150\pi^5\log2+56\pi^4G-270\pi^3\zeta(3)
               +1920\pi^2\beta(4)-675\pi\zeta(5).
\end{align}
Thus every coefficient and sign in the original five conjectural rows is
confirmed by exact finite algebra. This does not assert independence of
the remaining real components or of any single-value basis.

\subsection{The complete next even-weight table}

The same rule proves the following seven identities at weight eight.
These are extensions of the snapshot's displayed table; their global
bibliographic novelty is not asserted.
\begin{align}
 98304g_{7,1}={}&294912\beta(8)-320\pi^5\zeta(3)
                    -3072\pi^3\zeta(5)-24765\pi\zeta(7),\\
 49152g_{6,2}={}&-1024\pi^2\beta(6)-540672\beta(8)
        +320\pi^5\zeta(3)+6144\pi^3\zeta(5)+74295\pi\zeta(7),\\
 98304g_{5,3}={}&10240\pi^2\beta(6)+1671168\beta(8)
                   -18522\pi^3\zeta(5)-371475\pi\zeta(7),\\
 368640g_{4,4}={}&-224\pi^4\beta(4)-76800\pi^2\beta(6)
          -6635520\beta(8)+47430\pi^3\zeta(5)+1857375\pi\zeta(7),\\
 491520g_{3,5}={}&896\pi^4\beta(4)+102400\pi^2\beta(6)
                    +4915200\beta(8)\notag\\
          &-150\pi^5\zeta(3)-2700\pi^3\zeta(5)-1857375\pi\zeta(7),\\
 15482880g_{2,6}={}&-248\pi^6G-28224\pi^4\beta(4)
                  -1612800\pi^2\beta(6)-61931520\beta(8)\notag\\
          &+9450\pi^5\zeta(3)+56700\pi^3\zeta(5)+23402925\pi\zeta(7),\\
 30965760g_{1,7}={}&496\pi^6G+18816\pi^4\beta(4)+645120\pi^2\beta(6)
                    -5124\pi^7\log2\notag\\
          &-9450\pi^5\zeta(3)-28350\pi^3\zeta(5)-59535\pi\zeta(7).
\end{align}

\subsection{Independent differential proof and exact boundary data}

This provides a second derivation, with different bookkeeping, of every
coefficient in the two tables. Set $Q_n(L)=-\mathcal B_n(e^L)$, so
$Q_0=-1$ and $Q'_n=Q_{n-1}$ for $n\ge1$. Classical inversion is
\[
 \operatorname{Li}_n(z)+(-1)^n\operatorname{Li}_n(1/z)=Q_n(L).
\]
For $P_{a,b}(z)=D_{a,b}(z)-(-1)^{a+b}D_{a,b}(1/z)$ and $D=d/dL$,
termwise differentiation inside the disk followed by analytic continuation
gives
\[
 DP_{a,b}=P_{a-1,b}\quad(a>1),\qquad
 DP_{1,b}=-\operatorname{Li}_b(z)+\frac{Q_b(L)}{1-z}.
\]
Define the explicit function
\[
 R_{a,b}(L)=Q_{a+b}(L)-\operatorname{Li}_{a+b}(e^L)
 +\sum_{k=0}^{b}(-1)^k\binom{a+k-1}{k}
       Q_{b-k}(L)\operatorname{Li}_{a+k}(e^L).
\]
Pascal's identity gives $DR_{a,b}=R_{a-1,b}$ for $a>1$. At $a=1$,
the sum telescopes to
$DR_{1,b}=-\operatorname{Li}_b(z)+Q_b(L)/(1-z)$, with the $k=b$ term
$(-1)^bQ_0\operatorname{Li}_{b+1}$ essential. Hence
\begin{equation}\label{eq:gaussian-ode-certificate}
 P_{a,b}(z)=R_{a,b}(L)+\sum_{j=0}^{a-1}\frac{c_{a-j,b}}{j!}L^j,
\end{equation}
where
\begin{align}
 c_{1,b}&=(-1)^{b+1}b\zeta(b+1),\label{eq:gaussian-base-constant}\\
 c_{r,b}&=(1-(-1)^{r+b})\zeta(r,b)-R_{r,b}(0)\qquad(r\ge2).
 \label{eq:gaussian-other-constants}
\end{align}
In the second line every single polylogarithm in $R_{r,b}(0)$ is replaced
by its zeta value. The double-zeta convention is again descending:
$\zeta(r,b)=\sum_{n>m}n^{-r}m^{-b}$.

Here is a careful verification of the exceptional boundary line
\eqref{eq:gaussian-base-constant}. For $b>1$,
\[
 D_{1,b}(z)=\zeta(b)\operatorname{Li}_1(z)-\zeta(b,1)-\zeta(b+1)+o(1)
 \quad(z\to1),
\]
because the coefficient difference is absolutely summable and its value
at one is
$-\sum_{m\ge n\ge1}n^{-1}m^{-b}=-\zeta(b,1)-\zeta(b+1)$.
With $z=e^{i\theta}$, $\theta\downarrow0$,
$\operatorname{Li}_1(z)=-\log\theta+i\pi/2+o(1)$.
Thus for odd $b>1$, $P_{1,b}\to i\pi\zeta(b)$; for even $b$,
\[
 P_{1,b}=-2\zeta(b)\log\theta-2\zeta(b,1)-2\zeta(b+1)+o(1).
\]
Now use $Q_{2m}(0)=2\zeta(2m)$, $Q_1(0)=i\pi$,
$Q_{2m+1}(0)=0$ for $m\ge1$, and the two classical identities
\begin{align*}
 \zeta(b,1)&=\frac b2\zeta(b+1)
 -\frac12\sum_{r=2}^{b-1}\zeta(r)\zeta(b+1-r),\\
 \sum_{j=1}^{m-1}\zeta(2j)\zeta(2m-2j)&=\frac{2m+1}{2}\zeta(2m),\qquad m\ge2.
\end{align*}
They give exactly $P_{1,b}-R_{1,b}\to(-1)^{b+1}b\zeta(b+1)$.
For $b=1$ use $D_{1,1}(z)=\tfrac12\log^2(1-z)$ directly and
\[
 \Log(1-1/z)=\Log(1-z)-L+i\pi\quad(\Im z>0).
\]
This yields $P_{1,1}=R_{1,1}+\zeta(2)$, verifying the same formula.
For $r\ge2$, absolute convergence at $z=1$ proves
\eqref{eq:gaussian-other-constants}; induction then proves
\eqref{eq:gaussian-ode-certificate}.

\paragraph{Finite double-zeta certificates.}
Only odd-weight double zeta values at weights three, five, and seven are
needed for the two tables. They need no unproved reduction. The elementary
depth-two sum theorem is
\[
 \sum_{a=2}^{N-1}\zeta(a,N-a)=\zeta(N),\qquad N\ge3.
\]
Indeed, with $n=m+k$, the inner geometric sum is
\[
 \sum_{a=2}^{N-1}\frac1{n^a m^{N-a}}
 =\frac1{m^{N-1}}\left(\frac1k-\frac1{m+k}\right)
      -\frac1{k(m+k)^{N-1}}.
\]
Summing first over $k$ changes the first term to
$\sum_mH_m/m^{N-1}$, whereas the second is
$\sum_mH_{m-1}/m^{N-1}$. Their difference is $\zeta(N)$.
Summing stuffles then also proves the displayed height-one formula for
$\zeta(b,1)$. For completeness, the even-zeta convolution follows by
squaring $\sum_{j\ge1}\zeta(2j)x^{2j}=(1-\pi x\cot\pi x)/2$
and differentiating this elementary generating function.

At weight five, putting $X=\zeta(2,3)$, $Y=\zeta(3,2)$,
$Z=\zeta(4,1)$, the sum theorem, stuffle, and shuffle give
\[
 X+Y+Z=\zeta(5),\quad X+Y=\zeta(2)\zeta(3)-\zeta(5),\quad
 X+3Y+6Z=\zeta(2)\zeta(3).
\]
Their unique solution is
\[
 X=\tfrac92\zeta(5)-2\zeta(2)\zeta(3),\quad
 Y=3\zeta(2)\zeta(3)-\tfrac{11}2\zeta(5),\quad
 Z=2\zeta(5)-\zeta(2)\zeta(3).
\]
At weight seven let $x_a=\zeta(a,7-a)$ for $2\le a\le6$, and put
$A=\zeta(2)\zeta(5)$, $B=\zeta(3)\zeta(4)$, $Z=\zeta(7)$.
The corresponding exact system is
\[
 \begin{pmatrix}
 1&1&1&1&1\\1&0&0&1&0\\0&1&1&0&0\\
 1&2&3&5&10\\0&1&4&10&20
 \end{pmatrix}
 \begin{pmatrix}x_2\\x_3\\x_4\\x_5\\x_6\end{pmatrix}
 =\begin{pmatrix}Z\\A-Z\\B-Z\\A\\B\end{pmatrix}.
\]
Its determinant is $-2$. Thus
\[
 (x_2,x_3,x_4,x_5,x_6)
 =(10Z-4A-2B,\,-18Z+10A+B,\,17Z-10A,\,-11Z+5A+2B,\,3Z-A-B).
\]
The shuffle coefficients used here are exactly
\[
 \zeta(p)\zeta(q)=
 \sum_{j=0}^{p-1}\binom{q+j-1}{j}\zeta(q+j,p-j)
 +\sum_{j=0}^{q-1}\binom{p+j-1}{j}\zeta(p+j,q-j),\qquad p,q\ge2.
\]
Consequently both derivations of all twelve rows reduce to explicitly
displayed finite rational operations.

\paragraph{Reproducibility and validation.}
The script \texttt{gaussian\_parity.py} implements both
\eqref{eq:gaussian-inversion} and \eqref{eq:gaussian-real-coefficients}.
The independently organized script \texttt{extend\_weight8.py} implements
the differential certificate and solves the exact double-zeta systems.
All twelve identities agree symbolically. A separate 90-digit numerical
check evaluates
\[
 D_{a,b}(i)=\frac{i}{(a-1)!}\int_0^1
       \frac{(-\log t)^{a-1}\operatorname{Li}_b(it)}{1-it}\,dt,
\]
whose power-series expansion proves the formula for the integral directly.
The largest observed absolute residual over the twelve rows was below
$2.3\times10^{-90}$. These quadratures are numerical cross-checks; the
finite formulas and proofs establish the identities. Exact rational
coefficient vectors are recorded in \texttt{gaussian\_coefficients.json}.

